\element{fonctions}{
\begin{question}{fct-image-affine-1}
Soit $f(x)=1x+(-4)$. Calculer $f(1)$.
\begin{choiceshoriz}
\correctchoice{$-3$}
\wrongchoice{$-2$}
\wrongchoice{$5$}
\wrongchoice{$1$}
\end{choiceshoriz}
\end{question}
}

\element{fonctions}{
\begin{question}{fct-image-affine-2}
Soit $f(x)=2x+(-3)$. Calculer $f(2)$.
\begin{choiceshoriz}
\correctchoice{$1$}
\wrongchoice{$7$}
\wrongchoice{$-1$}
\wrongchoice{$2$}
\end{choiceshoriz}
\end{question}
}

\element{fonctions}{
\begin{question}{fct-image-affine-3}
Soit $f(x)=3x+(-2)$. Calculer $f(3)$.
\begin{choiceshoriz}
\correctchoice{$7$}
\wrongchoice{$4$}
\wrongchoice{$11$}
\wrongchoice{$1$}
\end{choiceshoriz}
\end{question}
}

\element{fonctions}{
\begin{question}{fct-image-affine-4}
Soit $f(x)=4x+(-1)$. Calculer $f(4)$.
\begin{choiceshoriz}
\correctchoice{$15$}
\wrongchoice{$7$}
\wrongchoice{$17$}
\wrongchoice{$3$}
\end{choiceshoriz}
\end{question}
}

\element{fonctions}{
\begin{question}{fct-image-affine-5}
Soit $f(x)=5x+(0)$. Calculer $f(5)$.
\begin{choiceshoriz}
\correctchoice{$25$}
\wrongchoice{$10$}
\wrongchoice{$5$}
\wrongchoice{$1$}
\end{choiceshoriz}
\end{question}
}

\element{fonctions}{
\begin{question}{fct-image-affine-6}
Soit $f(x)=1x+(1)$. Calculer $f(6)$.
\begin{choiceshoriz}
\correctchoice{$7$}
\wrongchoice{$8$}
\wrongchoice{$5$}
\wrongchoice{$1$}
\end{choiceshoriz}
\end{question}
}

\element{fonctions}{
\begin{question}{fct-image-affine-7}
Soit $f(x)=2x+(2)$. Calculer $f(7)$.
\begin{choiceshoriz}
\correctchoice{$16$}
\wrongchoice{$11$}
\wrongchoice{$12$}
\wrongchoice{$9$}
\end{choiceshoriz}
\end{question}
}

\element{fonctions}{
\begin{question}{fct-image-affine-8}
Soit $f(x)=3x+(3)$. Calculer $f(8)$.
\begin{choiceshoriz}
\correctchoice{$27$}
\wrongchoice{$14$}
\wrongchoice{$21$}
\wrongchoice{$11$}
\end{choiceshoriz}
\end{question}
}

\element{fonctions}{
\begin{question}{fct-image-affine-9}
Soit $f(x)=4x+(4)$. Calculer $f(9)$.
\begin{choiceshoriz}
\correctchoice{$40$}
\wrongchoice{$17$}
\wrongchoice{$32$}
\wrongchoice{$13$}
\end{choiceshoriz}
\end{question}
}

\element{fonctions}{
\begin{question}{fct-image-affine-10}
Soit $f(x)=5x+(5)$. Calculer $f(10)$.
\begin{choiceshoriz}
\correctchoice{$55$}
\wrongchoice{$20$}
\wrongchoice{$45$}
\wrongchoice{$15$}
\end{choiceshoriz}
\end{question}
}

\element{fonctions}{
\begin{question}{fct-antecedent-affine-1}
Soit $f(x)=2x+1$. Quel antécédent de $5$ ?
\begin{choiceshoriz}
\correctchoice{$2$}
\wrongchoice{$3$}
\wrongchoice{$4$}
\wrongchoice{$7$}
\end{choiceshoriz}
\end{question}
}

\element{fonctions}{
\begin{question}{fct-antecedent-affine-2}
Soit $f(x)=3x+2$. Quel antécédent de $11$ ?
\begin{choiceshoriz}
\correctchoice{$3$}
\wrongchoice{$4$}
\wrongchoice{$9$}
\wrongchoice{$14$}
\end{choiceshoriz}
\end{question}
}

\element{fonctions}{
\begin{question}{fct-antecedent-affine-3}
Soit $f(x)=4x+3$. Quel antécédent de $19$ ?
\begin{choiceshoriz}
\correctchoice{$4$}
\wrongchoice{$5$}
\wrongchoice{$16$}
\wrongchoice{$23$}
\end{choiceshoriz}
\end{question}
}

\element{fonctions}{
\begin{question}{fct-antecedent-affine-4}
Soit $f(x)=5x+4$. Quel antécédent de $29$ ?
\begin{choiceshoriz}
\correctchoice{$5$}
\wrongchoice{$6$}
\wrongchoice{$25$}
\wrongchoice{$34$}
\end{choiceshoriz}
\end{question}
}

\element{fonctions}{
\begin{question}{fct-antecedent-affine-5}
Soit $f(x)=6x+5$. Quel antécédent de $41$ ?
\begin{choiceshoriz}
\correctchoice{$6$}
\wrongchoice{$7$}
\wrongchoice{$36$}
\wrongchoice{$47$}
\end{choiceshoriz}
\end{question}
}

\element{fonctions}{
\begin{question}{fct-antecedent-affine-6}
Soit $f(x)=2x+6$. Quel antécédent de $20$ ?
\begin{choiceshoriz}
\correctchoice{$7$}
\wrongchoice{$8$}
\wrongchoice{$14$}
\wrongchoice{$22$}
\end{choiceshoriz}
\end{question}
}

\element{fonctions}{
\begin{question}{fct-antecedent-affine-7}
Soit $f(x)=3x+7$. Quel antécédent de $31$ ?
\begin{choiceshoriz}
\correctchoice{$8$}
\wrongchoice{$9$}
\wrongchoice{$24$}
\wrongchoice{$34$}
\end{choiceshoriz}
\end{question}
}

\element{fonctions}{
\begin{question}{fct-antecedent-affine-8}
Soit $f(x)=4x+8$. Quel antécédent de $44$ ?
\begin{choiceshoriz}
\correctchoice{$9$}
\wrongchoice{$10$}
\wrongchoice{$36$}
\wrongchoice{$48$}
\end{choiceshoriz}
\end{question}
}

\element{fonctions}{
\begin{question}{fct-antecedent-affine-9}
Soit $f(x)=5x+9$. Quel antécédent de $59$ ?
\begin{choiceshoriz}
\correctchoice{$10$}
\wrongchoice{$11$}
\wrongchoice{$50$}
\wrongchoice{$64$}
\end{choiceshoriz}
\end{question}
}

\element{fonctions}{
\begin{question}{fct-antecedent-affine-10}
Soit $f(x)=6x+10$. Quel antécédent de $76$ ?
\begin{choiceshoriz}
\correctchoice{$11$}
\wrongchoice{$12$}
\wrongchoice{$66$}
\wrongchoice{$82$}
\end{choiceshoriz}
\end{question}
}

\element{fonctions}{
\begin{question}{fct-lineaire-point-1}
La fonction linéaire $f(x)=-2x$ a pour représentation une droite passant par quel point ?
\begin{choiceshoriz}
\correctchoice{$(1;-2)$}
\wrongchoice{$(1;-1)$}
\wrongchoice{$(-2;1)$}
\wrongchoice{$1$}
\end{choiceshoriz}
\end{question}
}

\element{fonctions}{
\begin{question}{fct-lineaire-point-2}
La fonction linéaire $f(x)=-1x$ a pour représentation une droite passant par quel point ?
\begin{choiceshoriz}
\correctchoice{$(2;-2)$}
\wrongchoice{$(2;-1)$}
\wrongchoice{$(-2;2)$}
\wrongchoice{$(2;1)$}
\end{choiceshoriz}
\end{question}
}

\element{fonctions}{
\begin{question}{fct-lineaire-point-3}
La fonction linéaire $f(x)=0x$ a pour représentation une droite passant par quel point ?
\begin{choiceshoriz}
\correctchoice{$(3;0)$}
\wrongchoice{$(3;1)$}
\wrongchoice{$(0;3)$}
\wrongchoice{$(3;3)$}
\end{choiceshoriz}
\end{question}
}

\element{fonctions}{
\begin{question}{fct-lineaire-point-4}
La fonction linéaire $f(x)=1x$ a pour représentation une droite passant par quel point ?
\begin{choiceshoriz}
\correctchoice{$(4;4)$}
\wrongchoice{$(4;5)$}
\wrongchoice{$1$}
\wrongchoice{$2$}
\end{choiceshoriz}
\end{question}
}

\element{fonctions}{
\begin{question}{fct-lineaire-point-5}
La fonction linéaire $f(x)=2x$ a pour représentation une droite passant par quel point ?
\begin{choiceshoriz}
\correctchoice{$(5;10)$}
\wrongchoice{$(5;11)$}
\wrongchoice{$(10;5)$}
\wrongchoice{$(5;7)$}
\end{choiceshoriz}
\end{question}
}

\element{fonctions}{
\begin{question}{fct-lineaire-point-6}
La fonction linéaire $f(x)=3x$ a pour représentation une droite passant par quel point ?
\begin{choiceshoriz}
\correctchoice{$(6;18)$}
\wrongchoice{$(6;19)$}
\wrongchoice{$(18;6)$}
\wrongchoice{$(6;9)$}
\end{choiceshoriz}
\end{question}
}

\element{fonctions}{
\begin{question}{fct-lineaire-point-7}
La fonction linéaire $f(x)=-2x$ a pour représentation une droite passant par quel point ?
\begin{choiceshoriz}
\correctchoice{$(7;-14)$}
\wrongchoice{$(7;-13)$}
\wrongchoice{$(-14;7)$}
\wrongchoice{$(7;5)$}
\end{choiceshoriz}
\end{question}
}

\element{fonctions}{
\begin{question}{fct-lineaire-point-8}
La fonction linéaire $f(x)=-1x$ a pour représentation une droite passant par quel point ?
\begin{choiceshoriz}
\correctchoice{$(8;-8)$}
\wrongchoice{$(8;-7)$}
\wrongchoice{$(-8;8)$}
\wrongchoice{$(8;7)$}
\end{choiceshoriz}
\end{question}
}

\element{fonctions}{
\begin{question}{fct-lineaire-point-9}
La fonction linéaire $f(x)=0x$ a pour représentation une droite passant par quel point ?
\begin{choiceshoriz}
\correctchoice{$(9;0)$}
\wrongchoice{$(9;1)$}
\wrongchoice{$(0;9)$}
\wrongchoice{$(9;9)$}
\end{choiceshoriz}
\end{question}
}

\element{fonctions}{
\begin{question}{fct-lineaire-point-10}
La fonction linéaire $f(x)=1x$ a pour représentation une droite passant par quel point ?
\begin{choiceshoriz}
\correctchoice{$(10;10)$}
\wrongchoice{$(10;11)$}
\wrongchoice{$1$}
\wrongchoice{$2$}
\end{choiceshoriz}
\end{question}
}

\element{fonctions}{
\begin{question}{fct-coeff-directeur-1}
Une droite passe par $(0;2)$ et $(2;6)$. Son coefficient directeur est :
\begin{choiceshoriz}
\correctchoice{$2$}
\wrongchoice{$4$}
\wrongchoice{$1$}
\wrongchoice{$3$}
\end{choiceshoriz}
\end{question}
}

\element{fonctions}{
\begin{question}{fct-coeff-directeur-2}
Une droite passe par $(0;3)$ et $(2;7)$. Son coefficient directeur est :
\begin{choiceshoriz}
\correctchoice{$2$}
\wrongchoice{$4$}
\wrongchoice{$5$}
\wrongchoice{$3$}
\end{choiceshoriz}
\end{question}
}

\element{fonctions}{
\begin{question}{fct-coeff-directeur-3}
Une droite passe par $(0;4)$ et $(2;8)$. Son coefficient directeur est :
\begin{choiceshoriz}
\correctchoice{$2$}
\wrongchoice{$4$}
\wrongchoice{$6$}
\wrongchoice{$1$}
\end{choiceshoriz}
\end{question}
}

\element{fonctions}{
\begin{question}{fct-coeff-directeur-4}
Une droite passe par $(0;5)$ et $(2;9)$. Son coefficient directeur est :
\begin{choiceshoriz}
\correctchoice{$2$}
\wrongchoice{$4$}
\wrongchoice{$7$}
\wrongchoice{$5$}
\end{choiceshoriz}
\end{question}
}

\element{fonctions}{
\begin{question}{fct-coeff-directeur-5}
Une droite passe par $(0;6)$ et $(2;10)$. Son coefficient directeur est :
\begin{choiceshoriz}
\correctchoice{$2$}
\wrongchoice{$4$}
\wrongchoice{$8$}
\wrongchoice{$6$}
\end{choiceshoriz}
\end{question}
}

\element{fonctions}{
\begin{question}{fct-coeff-directeur-6}
Une droite passe par $(0;7)$ et $(2;11)$. Son coefficient directeur est :
\begin{choiceshoriz}
\correctchoice{$2$}
\wrongchoice{$4$}
\wrongchoice{$9$}
\wrongchoice{$7$}
\end{choiceshoriz}
\end{question}
}

\element{fonctions}{
\begin{question}{fct-coeff-directeur-7}
Une droite passe par $(0;8)$ et $(2;12)$. Son coefficient directeur est :
\begin{choiceshoriz}
\correctchoice{$2$}
\wrongchoice{$4$}
\wrongchoice{$10$}
\wrongchoice{$8$}
\end{choiceshoriz}
\end{question}
}

\element{fonctions}{
\begin{question}{fct-coeff-directeur-8}
Une droite passe par $(0;9)$ et $(2;13)$. Son coefficient directeur est :
\begin{choiceshoriz}
\correctchoice{$2$}
\wrongchoice{$4$}
\wrongchoice{$11$}
\wrongchoice{$9$}
\end{choiceshoriz}
\end{question}
}

\element{fonctions}{
\begin{question}{fct-coeff-directeur-9}
Une droite passe par $(0;10)$ et $(2;14)$. Son coefficient directeur est :
\begin{choiceshoriz}
\correctchoice{$2$}
\wrongchoice{$4$}
\wrongchoice{$12$}
\wrongchoice{$10$}
\end{choiceshoriz}
\end{question}
}

\element{fonctions}{
\begin{question}{fct-coeff-directeur-10}
Une droite passe par $(0;11)$ et $(2;15)$. Son coefficient directeur est :
\begin{choiceshoriz}
\correctchoice{$2$}
\wrongchoice{$4$}
\wrongchoice{$13$}
\wrongchoice{$11$}
\end{choiceshoriz}
\end{question}
}

\element{fonctions}{
\begin{question}{fct-ordonnee-origine-1}
Pour $f(x)=1x+(-3)$, l'ordonnée à l'origine est :
\begin{choiceshoriz}
\correctchoice{$-3$}
\wrongchoice{$1$}
\wrongchoice{$-2$}
\wrongchoice{$3$}
\end{choiceshoriz}
\end{question}
}

\element{fonctions}{
\begin{question}{fct-ordonnee-origine-2}
Pour $f(x)=2x+(-2)$, l'ordonnée à l'origine est :
\begin{choiceshoriz}
\correctchoice{$-2$}
\wrongchoice{$2$}
\wrongchoice{$0$}
\wrongchoice{$1$}
\end{choiceshoriz}
\end{question}
}

\element{fonctions}{
\begin{question}{fct-ordonnee-origine-3}
Pour $f(x)=3x+(-1)$, l'ordonnée à l'origine est :
\begin{choiceshoriz}
\correctchoice{$-1$}
\wrongchoice{$3$}
\wrongchoice{$2$}
\wrongchoice{$1$}
\end{choiceshoriz}
\end{question}
}

\element{fonctions}{
\begin{question}{fct-ordonnee-origine-4}
Pour $f(x)=4x+(0)$, l'ordonnée à l'origine est :
\begin{choiceshoriz}
\correctchoice{$0$}
\wrongchoice{$4$}
\wrongchoice{$1$}
\wrongchoice{$2$}
\end{choiceshoriz}
\end{question}
}

\element{fonctions}{
\begin{question}{fct-ordonnee-origine-5}
Pour $f(x)=1x+(1)$, l'ordonnée à l'origine est :
\begin{choiceshoriz}
\correctchoice{$1$}
\wrongchoice{$2$}
\wrongchoice{$-1$}
\wrongchoice{$3$}
\end{choiceshoriz}
\end{question}
}

\element{fonctions}{
\begin{question}{fct-ordonnee-origine-6}
Pour $f(x)=2x+(2)$, l'ordonnée à l'origine est :
\begin{choiceshoriz}
\correctchoice{$2$}
\wrongchoice{$4$}
\wrongchoice{$-2$}
\wrongchoice{$1$}
\end{choiceshoriz}
\end{question}
}

\element{fonctions}{
\begin{question}{fct-ordonnee-origine-7}
Pour $f(x)=3x+(3)$, l'ordonnée à l'origine est :
\begin{choiceshoriz}
\correctchoice{$3$}
\wrongchoice{$6$}
\wrongchoice{$-3$}
\wrongchoice{$1$}
\end{choiceshoriz}
\end{question}
}

\element{fonctions}{
\begin{question}{fct-ordonnee-origine-8}
Pour $f(x)=4x+(4)$, l'ordonnée à l'origine est :
\begin{choiceshoriz}
\correctchoice{$4$}
\wrongchoice{$8$}
\wrongchoice{$-4$}
\wrongchoice{$1$}
\end{choiceshoriz}
\end{question}
}

\element{fonctions}{
\begin{question}{fct-ordonnee-origine-9}
Pour $f(x)=1x+(5)$, l'ordonnée à l'origine est :
\begin{choiceshoriz}
\correctchoice{$5$}
\wrongchoice{$1$}
\wrongchoice{$6$}
\wrongchoice{$-5$}
\end{choiceshoriz}
\end{question}
}

\element{fonctions}{
\begin{question}{fct-ordonnee-origine-10}
Pour $f(x)=2x+(6)$, l'ordonnée à l'origine est :
\begin{choiceshoriz}
\correctchoice{$6$}
\wrongchoice{$2$}
\wrongchoice{$8$}
\wrongchoice{$-6$}
\end{choiceshoriz}
\end{question}
}

\element{fonctions}{
\begin{question}{fct-variation-lineaire-1}
La fonction $f(x)=-2x+1$ est :
\begin{choiceshoriz}
\correctchoice{décroissante}
\wrongchoice{croissante}
\wrongchoice{constante}
\wrongchoice{ni affine ni linéaire}
\end{choiceshoriz}
\end{question}
}

\element{fonctions}{
\begin{question}{fct-variation-lineaire-2}
La fonction $f(x)=-1x+1$ est :
\begin{choiceshoriz}
\correctchoice{décroissante}
\wrongchoice{croissante}
\wrongchoice{constante}
\wrongchoice{ni affine ni linéaire}
\end{choiceshoriz}
\end{question}
}

\element{fonctions}{
\begin{question}{fct-variation-lineaire-3}
La fonction $f(x)=0x+1$ est :
\begin{choiceshoriz}
\correctchoice{constante}
\wrongchoice{croissante}
\wrongchoice{décroissante}
\wrongchoice{ni affine ni linéaire}
\end{choiceshoriz}
\end{question}
}

\element{fonctions}{
\begin{question}{fct-variation-lineaire-4}
La fonction $f(x)=1x+1$ est :
\begin{choiceshoriz}
\correctchoice{croissante}
\wrongchoice{décroissante}
\wrongchoice{constante}
\wrongchoice{ni affine ni linéaire}
\end{choiceshoriz}
\end{question}
}

\element{fonctions}{
\begin{question}{fct-variation-lineaire-5}
La fonction $f(x)=2x+1$ est :
\begin{choiceshoriz}
\correctchoice{croissante}
\wrongchoice{décroissante}
\wrongchoice{constante}
\wrongchoice{ni affine ni linéaire}
\end{choiceshoriz}
\end{question}
}

\element{fonctions}{
\begin{question}{fct-variation-lineaire-6}
La fonction $f(x)=-2x+1$ est :
\begin{choiceshoriz}
\correctchoice{décroissante}
\wrongchoice{croissante}
\wrongchoice{constante}
\wrongchoice{ni affine ni linéaire}
\end{choiceshoriz}
\end{question}
}

\element{fonctions}{
\begin{question}{fct-variation-lineaire-7}
La fonction $f(x)=-1x+1$ est :
\begin{choiceshoriz}
\correctchoice{décroissante}
\wrongchoice{croissante}
\wrongchoice{constante}
\wrongchoice{ni affine ni linéaire}
\end{choiceshoriz}
\end{question}
}

\element{fonctions}{
\begin{question}{fct-variation-lineaire-8}
La fonction $f(x)=0x+1$ est :
\begin{choiceshoriz}
\correctchoice{constante}
\wrongchoice{croissante}
\wrongchoice{décroissante}
\wrongchoice{ni affine ni linéaire}
\end{choiceshoriz}
\end{question}
}

\element{fonctions}{
\begin{question}{fct-variation-lineaire-9}
La fonction $f(x)=1x+1$ est :
\begin{choiceshoriz}
\correctchoice{croissante}
\wrongchoice{décroissante}
\wrongchoice{constante}
\wrongchoice{ni affine ni linéaire}
\end{choiceshoriz}
\end{question}
}

\element{fonctions}{
\begin{question}{fct-variation-lineaire-10}
La fonction $f(x)=2x+1$ est :
\begin{choiceshoriz}
\correctchoice{croissante}
\wrongchoice{décroissante}
\wrongchoice{constante}
\wrongchoice{ni affine ni linéaire}
\end{choiceshoriz}
\end{question}
}

\element{fonctions}{
\begin{question}{fct-appartenance-1}
Le point $(3;5)$ appartient à la courbe de quelle fonction ?
\begin{choiceshoriz}
\correctchoice{$f(x)=1x+2$}
\wrongchoice{$f(x)=2x+1$}
\wrongchoice{$f(x)=1x-2$}
\wrongchoice{$f(x)=x+5$}
\end{choiceshoriz}
\end{question}
}

\element{fonctions}{
\begin{question}{fct-appartenance-2}
Le point $(4;11)$ appartient à la courbe de quelle fonction ?
\begin{choiceshoriz}
\correctchoice{$f(x)=2x+3$}
\wrongchoice{$f(x)=3x+2$}
\wrongchoice{$f(x)=2x-3$}
\wrongchoice{$f(x)=x+11$}
\end{choiceshoriz}
\end{question}
}

\element{fonctions}{
\begin{question}{fct-appartenance-3}
Le point $(5;19)$ appartient à la courbe de quelle fonction ?
\begin{choiceshoriz}
\correctchoice{$f(x)=3x+4$}
\wrongchoice{$f(x)=4x+3$}
\wrongchoice{$f(x)=3x-4$}
\wrongchoice{$f(x)=x+19$}
\end{choiceshoriz}
\end{question}
}

\element{fonctions}{
\begin{question}{fct-appartenance-4}
Le point $(6;29)$ appartient à la courbe de quelle fonction ?
\begin{choiceshoriz}
\correctchoice{$f(x)=4x+5$}
\wrongchoice{$f(x)=5x+4$}
\wrongchoice{$f(x)=4x-5$}
\wrongchoice{$f(x)=x+29$}
\end{choiceshoriz}
\end{question}
}

\element{fonctions}{
\begin{question}{fct-appartenance-5}
Le point $(7;41)$ appartient à la courbe de quelle fonction ?
\begin{choiceshoriz}
\correctchoice{$f(x)=5x+6$}
\wrongchoice{$f(x)=6x+5$}
\wrongchoice{$f(x)=5x-6$}
\wrongchoice{$f(x)=x+41$}
\end{choiceshoriz}
\end{question}
}

\element{fonctions}{
\begin{question}{fct-appartenance-6}
Le point $(8;15)$ appartient à la courbe de quelle fonction ?
\begin{choiceshoriz}
\correctchoice{$f(x)=1x+7$}
\wrongchoice{$f(x)=7x+1$}
\wrongchoice{$f(x)=1x-7$}
\wrongchoice{$f(x)=x+15$}
\end{choiceshoriz}
\end{question}
}

\element{fonctions}{
\begin{question}{fct-appartenance-7}
Le point $(9;26)$ appartient à la courbe de quelle fonction ?
\begin{choiceshoriz}
\correctchoice{$f(x)=2x+8$}
\wrongchoice{$f(x)=8x+2$}
\wrongchoice{$f(x)=2x-8$}
\wrongchoice{$f(x)=x+26$}
\end{choiceshoriz}
\end{question}
}

\element{fonctions}{
\begin{question}{fct-appartenance-8}
Le point $(10;39)$ appartient à la courbe de quelle fonction ?
\begin{choiceshoriz}
\correctchoice{$f(x)=3x+9$}
\wrongchoice{$f(x)=9x+3$}
\wrongchoice{$f(x)=3x-9$}
\wrongchoice{$f(x)=x+39$}
\end{choiceshoriz}
\end{question}
}

\element{fonctions}{
\begin{question}{fct-appartenance-9}
Le point $(11;54)$ appartient à la courbe de quelle fonction ?
\begin{choiceshoriz}
\correctchoice{$f(x)=4x+10$}
\wrongchoice{$f(x)=10x+4$}
\wrongchoice{$f(x)=4x-10$}
\wrongchoice{$f(x)=x+54$}
\end{choiceshoriz}
\end{question}
}

\element{fonctions}{
\begin{question}{fct-appartenance-10}
Le point $(12;71)$ appartient à la courbe de quelle fonction ?
\begin{choiceshoriz}
\correctchoice{$f(x)=5x+11$}
\wrongchoice{$f(x)=11x+5$}
\wrongchoice{$f(x)=5x-11$}
\wrongchoice{$f(x)=x+71$}
\end{choiceshoriz}
\end{question}
}
